\documentclass[11pt]{article}

\usepackage[margin=1in]{geometry}
\usepackage[T1]{fontenc}
\usepackage[utf8]{inputenc}
\usepackage{longtable}
\usepackage{booktabs}
\usepackage{array}
\usepackage{hyperref}
\usepackage{listings}
\usepackage{xcolor}

\hypersetup{
  colorlinks=true,
  linkcolor=blue,
  urlcolor=blue,
  citecolor=blue
}

\lstset{
  basicstyle=\ttfamily\small,
  breaklines=true,
  frame=single,
  columns=fullflexible
}

\title{RES201 Week 1 Report\\Chemical-Family Domain Shift in Pretrained ALIGNN}
\author{Prepared from workspace audit of \texttt{res201-alignn-domain-shift}}
\date{April 7, 2026}

\begin{document}
\maketitle

\section{Purpose}
This report documents the complete Week 1 execution status for the RES201 project
\textit{Chemical-Family Domain Shift in Pretrained ALIGNN Data Efficiency for Oxide vs. Nitride Formation Energy Prediction}.
It is intended to serve two purposes:
\begin{enumerate}
  \item provide a reproducible record of how the Week 1 deliverables were executed and verified; and
  \item provide reusable methods text and artifact references for the final research paper.
\end{enumerate}

The official Week 1 deliverables, taken from the project brief in
\texttt{Project\_Task/RES201 Project.pdf}, were:
\begin{itemize}
  \item oxide subset filtered and confirmed;
  \item nitride subset filtered and confirmed with oxynitrides excluded from the nitride arm;
  \item original JARVIS split IDs applied before family filtering;
  \item fixed oxide and nitride test sets;
  \item pretrained checkpoint loaded;
  \item oxide zero-shot MAE reported;
  \item nitride zero-shot MAE reported;
  \item one oxide fine-tuning run at $N=50$ completed;
  \item one nitride fine-tuning run at $N=50$ completed.
\end{itemize}

\section{High-Level Outcome}
Week 1 is complete in this workspace. Both family datasets exist and validate successfully, both zero-shot evaluations exist, and both $N=50$ fine-tuning runs exist with saved checkpoints, histories, test predictions, and summaries.

The most important scientific result from Week 1 is that under the current setup the $N=50$ fine-tune runs performed worse than zero-shot for both families. This does \textbf{not} invalidate Week 1 completion; instead it becomes a key observation to investigate in Week 2 and later analysis.

\section{Methodology}

\subsection{Stage 2 Dataset Construction}
The dataset source used in this workspace is \texttt{dft\_3d\_2021} from JARVIS.
The project rule from the brief was followed exactly:
\begin{enumerate}
  \item recover or provide the original benchmark train/validation/test split IDs;
  \item apply those split IDs first;
  \item only then filter by chemical family.
\end{enumerate}

The Stage 2 pipeline is implemented in:
\begin{itemize}
  \item \texttt{scripts/dataset/build\_res201\_family\_datasets.py}
  \item \texttt{scripts/dataset/res201\_stage2\_lib.py}
  \item \texttt{scripts/dataset/validate\_res201\_stage2.py}
\end{itemize}

The split manifest used for formation energy is:
\begin{itemize}
  \item \texttt{manifests/dft\_3d\_formation\_energy\_peratom\_splits.csv}
\end{itemize}
and the explicitly conflicting IDs are recorded in:
\begin{itemize}
  \item \texttt{manifests/dft\_3d\_formation\_energy\_peratom\_splits\_conflicts.csv}
\end{itemize}

The family definitions implemented in the workspace are:
\begin{itemize}
  \item oxide arm: any material containing O;
  \item nitride arm: materials containing N and not O;
  \item oxynitrides: excluded from the nitride arm;
  \item oxynitrides may still appear in the oxide arm because the brief defines oxides by O presence.
\end{itemize}

This choice is visible in the saved family summaries:
\begin{itemize}
  \item oxide summary reports \texttt{oxynitride\_count\_in\_family = 499};
  \item nitride summary reports \texttt{oxynitride\_count\_in\_family = 0}.
\end{itemize}

\subsection{Stage 2 Validation}
The validator was rerun directly on the current workspace:

\begin{lstlisting}[language=bash]
python scripts/dataset/validate_res201_stage2.py --root ./data_shared
\end{lstlisting}

It returned \texttt{status = ok}. The validated family counts are:
\begin{center}
\begin{tabular}{lrrrrr}
\toprule
Family & All & Train & Val & Test & Pool \\
\midrule
Oxide   & 14991 & 11960 & 1547 & 1484 & 13507 \\
Nitride &  2288 &  1837 &  209 &  242 &  2046 \\
\bottomrule
\end{tabular}
\end{center}

These values match the saved summaries in \texttt{data\_shared/*/summaries/summary.json}.

\subsection{Stage 3 Environment and Pretrained Checkpoint}
The Week 1 training environment is documented by:
\begin{itemize}
  \item \texttt{env/bootstrap\_res201\_stage3\_train.sh}
  \item \texttt{requirements/res201\_train\_frozen.txt}
\end{itemize}

The locally extracted pretrained formation-energy checkpoint used for Week 1 is:
\begin{itemize}
  \item \texttt{jv\_formation\_energy\_peratom\_alignn/checkpoint\_300.pt}
  \item \texttt{jv\_formation\_energy\_peratom\_alignn/config.json}
\end{itemize}

During the audit it was confirmed that the installed ALIGNN version in the
\texttt{res201\_train} environment differs in important ways from the assumptions in the earlier timeline notes and tutorial summary.
Because of those API mismatches, the Week 1 fine-tuning path was repaired locally so that the current environment can reproducibly:
\begin{itemize}
  \item load user-provided \texttt{id\_prop.csv} plus POSCAR data from a local dataset root;
  \item load the local pretrained ALIGNN checkpoint correctly;
  \item freeze all parameters except the chosen final groups;
  \item train on the fixed train/validation/test partition for the chosen $N$.
\end{itemize}

\subsection{Zero-Shot Evaluation Procedure}
The zero-shot evaluation script is:
\begin{itemize}
  \item \texttt{scripts/shared/zero\_shot\_alignn.py}
\end{itemize}

This script:
\begin{enumerate}
  \item reads the fixed family test manifest;
  \item loads each POSCAR structure from \texttt{data\_shared/<family>/structures};
  \item runs \texttt{alignn.pretrained.get\_prediction};
  \item coerces ALIGNN outputs to scalar floats;
  \item computes test MAE in eV/atom;
  \item writes a predictions CSV and updates the canonical report-level zero-shot summary table.
\end{enumerate}

Representative command form:
\begin{lstlisting}[language=bash]
source "$(conda info --base)/etc/profile.d/conda.sh"
conda activate res201_train
python scripts/shared/Evaluate_ALIGNN_Zero_Shot.py --family oxide --repo-root .
python scripts/shared/Evaluate_ALIGNN_Zero_Shot.py --family nitride --repo-root .
\end{lstlisting}

\subsection{$N=50$ Fine-Tuning Dataset Preparation}
The Week 1 fine-tuning dataset builder is:
\begin{itemize}
  \item \texttt{scripts/shared/prepare\_week1\_finetune\_dataset.py}
\end{itemize}

This script constructs a local \texttt{dataset\_root} for each family as follows:
\begin{enumerate}
  \item read the fixed family \texttt{pool.csv} and \texttt{test.csv};
  \item sample $N=50$ points from the pool using NumPy RNG with \texttt{seed = 0};
  \item allocate $45$ samples to train and $5$ samples to validation;
  \item keep the family test set fixed and unchanged;
  \item copy the corresponding POSCAR files into \texttt{Results_Before_Correction/<family>/N50\_seed0/dataset\_root};
  \item write \texttt{id\_prop.csv} and \texttt{split\_manifest.json}.
\end{enumerate}

Representative commands used:
\begin{lstlisting}[language=bash]
python scripts/shared/Prepare_Week1_Finetuning_Dataset.py \
  --family oxide --N 50 --seed 0 --repo-root . --link-mode copy

python scripts/shared/Prepare_Week1_Finetuning_Dataset.py \
  --family nitride --N 50 --seed 0 --repo-root . --link-mode copy
\end{lstlisting}

\subsection{$N=50$ Fine-Tuning Configuration}
The config writer is:
\begin{itemize}
  \item \texttt{scripts/shared/write\_week1\_alignn\_config.py}
\end{itemize}

It writes a TrainingConfig-compatible JSON for the currently installed ALIGNN version with:
\begin{itemize}
  \item \texttt{dataset = user\_data};
  \item \texttt{target = target};
  \item \texttt{filename = id\_prop.csv};
  \item \texttt{epochs = 50};
  \item \texttt{batch\_size = 16};
  \item \texttt{learning\_rate = 1e-4};
  \item fixed \texttt{n\_train}, \texttt{n\_val}, and \texttt{n\_test};
  \item the pretrained ALIGNN model architecture copied from the local checkpoint config.
\end{itemize}

The generated Week 1 configs are:
\begin{itemize}
  \item \texttt{configs/oxide\_week1\_N50\_seed0.finetune\_last2.json}
  \item \texttt{configs/nitride\_week1\_N50\_seed0.finetune\_last2.json}
\end{itemize}

\subsection{$N=50$ Fine-Tuning Training Procedure}
The training wrapper used for Week 1 is:
\begin{itemize}
  \item \texttt{scripts/shared/finetune\_last2\_alignn.py}
\end{itemize}

This wrapper was used because the installed ALIGNN CLI and lower-level APIs did not directly support the desired pretrained user-data fine-tuning path for this checkpoint version.
The Week 1 wrapper performs the following:
\begin{enumerate}
  \item load the local pretrained checkpoint from \texttt{checkpoint\_300.pt};
  \item rebuild the ALIGNN model from the saved checkpoint config;
  \item freeze all model parameters;
  \item unfreeze only the final two parameter groups used in this workspace:
        \texttt{gcn\_layers.3} and \texttt{fc};
  \item build LMDB-backed train/validation/test graph datasets from the local \texttt{dataset\_root};
  \item train for $50$ epochs using AdamW and one-cycle LR scheduling;
  \item save checkpoint, history, and test prediction artifacts.
\end{enumerate}

Representative commands used:
\begin{lstlisting}[language=bash]
python scripts/shared/Fine_Tune_Last_Two_ALIGNN_Layers.py \
  --config configs/oxide_week1_N50_seed0.finetune_last2.json \
  --output-dir Results_Before_Correction/oxide/N50_seed0/finetune_last2 \
  --dataset-root Results_Before_Correction/oxide/N50_seed0/dataset_root \
  --pretrained-checkpoint jv_formation_energy_peratom_alignn/checkpoint_300.pt \
  --pretrained-config jv_formation_energy_peratom_alignn/config.json \
  --device cpu

python scripts/shared/Fine_Tune_Last_Two_ALIGNN_Layers.py \
  --config configs/nitride_week1_N50_seed0.finetune_last2.json \
  --output-dir Results_Before_Correction/nitride/N50_seed0/finetune_last2 \
  --dataset-root Results_Before_Correction/nitride/N50_seed0/dataset_root \
  --pretrained-checkpoint jv_formation_energy_peratom_alignn/checkpoint_300.pt \
  --pretrained-config jv_formation_energy_peratom_alignn/config.json \
  --device cpu
\end{lstlisting}

\section{Week 1 Results}

\subsection{Zero-Shot Results}
\begin{center}
\begin{tabular}{lrr}
\toprule
Family & Test size & Zero-shot MAE (eV/atom) \\
\midrule
Oxide   & 1484 & 0.0341836068 \\
Nitride &  242 & 0.0695420150 \\
\bottomrule
\end{tabular}
\end{center}

\subsection{$N=50$ Fine-Tuning Results}
\begin{center}
\begin{tabular}{lrrrr}
\toprule
Family & Train & Val & Test & Test MAE (eV/atom) \\
\midrule
Oxide   & 45 & 5 & 1484 & 0.0700507342 \\
Nitride & 45 & 5 &  242 & 0.1617631432 \\
\bottomrule
\end{tabular}
\end{center}

\subsection{Interpretation}
For this first Week 1 pass, both zero-shot baselines outperform the corresponding $N=50$ fine-tuning runs.
That means Week 1 succeeded as an experimental milestone, but the numerical outcome suggests that:
\begin{itemize}
  \item the present fine-tuning protocol is not yet data-efficient at $N=50$;
  \item additional seeds, larger $N$, and from-scratch baselines remain essential;
  \item the current result should be reported honestly rather than hidden.
\end{itemize}

\section{Week 1 Deliverable Checklist}
\begin{longtable}{p{0.34\linewidth}p{0.10\linewidth}p{0.48\linewidth}}
\toprule
Deliverable & Status & Evidence \\
\midrule
Oxide subset filtered and confirmed & Done & \texttt{data\_shared/oxide/manifests/*.csv}, \texttt{data\_shared/oxide/summaries/summary.json} \\
Nitride subset filtered and confirmed & Done & \texttt{data\_shared/nitride/manifests/*.csv}, \texttt{data\_shared/nitride/summaries/summary.json} \\
Oxynitrides excluded from nitride arm & Done & nitride summary reports \texttt{oxynitride\_count\_in\_family = 0} \\
Original split IDs applied first & Done & split source in summaries points to \texttt{manifests/dft\_3d\_formation\_energy\_peratom\_splits.csv} \\
Fixed oxide test set & Done & \texttt{data\_shared/oxide/manifests/test.csv} \\
Fixed nitride test set & Done & \texttt{data\_shared/nitride/manifests/test.csv} \\
Pretrained checkpoint loaded & Done & \texttt{jv\_formation\_energy\_peratom\_alignn/checkpoint\_300.pt} \\
Oxide zero-shot MAE reported & Done & \texttt{reports/zero\_shot/zero\_shot\_summary.csv} \\
Nitride zero-shot MAE reported & Done & \texttt{reports/zero\_shot/zero\_shot\_summary.csv} \\
One oxide fine-tune at $N=50$ completed & Done & \texttt{Results_Before_Correction/oxide/N50\_seed0/finetune\_last2/summary.json} \\
One nitride fine-tune at $N=50$ completed & Done & \texttt{Results_Before_Correction/nitride/N50\_seed0/finetune\_last2/summary.json} \\
\bottomrule
\end{longtable}

\section{Verification Commands}
The simplest verification command is now:
\begin{lstlisting}[language=bash]
scripts/shared/Check_Week1_Baseline_Status.sh .
\end{lstlisting}

Additional direct checks:
\begin{lstlisting}[language=bash]
python scripts/dataset/validate_res201_stage2.py --root ./data_shared

python - <<'PY'
import csv
for row in csv.DictReader(open("reports/zero_shot/zero_shot_summary.csv")):
    print(row)
PY
python -m json.tool Results_Before_Correction/oxide/N50_seed0/finetune_last2/summary.json
python -m json.tool Results_Before_Correction/nitride/N50_seed0/finetune_last2/summary.json

python - <<'PY'
import json
for p in [
    "Results_Before_Correction/oxide/N50_seed0/finetune_last2/history_val.json",
    "Results_Before_Correction/nitride/N50_seed0/finetune_last2/history_val.json",
]:
    h = json.load(open(p))
    print(p, len(h), h[-1])
PY
\end{lstlisting}

\section{Workspace Inventory}

\subsection{Top-Level Files and Folders}
\begin{longtable}{p{0.28\linewidth}p{0.64\linewidth}}
\toprule
Path & Description \\
\midrule
\texttt{.codex} & Local Codex marker file in this workspace. \\
\texttt{.git} & Git metadata for the local repository clone. \\
\texttt{.gitignore} & Ignore rules controlling which generated artifacts should stay local. \\
\texttt{Project\_Task/} & Course and project reference PDFs: project brief, ALIGNN tutorial, and timeline notes. \\
\texttt{README.md} & Original Stage 2 payload README. \\
\texttt{configs/} & Saved Week 1 configuration JSON files, including historical failed attempts and final working fine-tune configs. \\
\texttt{data\_shared/} & Canonical Stage 2 family datasets, manifests, summaries, diagnostics, structures, and ALIGNN-ready exports. \\
\texttt{docs/} & Stage 2 and Stage 3 playbooks plus frozen dataset notes. \\
\texttt{env/} & Environment bootstrap and validation scripts for the data and training environments. \\
\texttt{jv\_formation\_energy\_peratom\_alignn/} & Local pretrained checkpoint directory containing the extracted formation-energy checkpoint and its config. \\
\texttt{manifests/} & Official split manifest, conflict list, and split manifest templates. \\
\texttt{notebooks/} & Currently empty notebook area reserved for later analysis. \\
\texttt{reports/} & Reporting area. This report and generated inventory files belong here. \\
\texttt{requirements/} & Frozen dependency record for the training environment. \\
\texttt{res201\_stage2\_data\_shared.zip} & Local archive copy of the Stage 2 shared dataset payload. \\
\texttt{Results_Before_Correction/} & Week 1 outputs, including zero-shot predictions and $N=50$ oxide/nitride experiment folders. \\
\texttt{scripts/} & Dataset pipeline scripts and shared Week 1 experiment scripts. \\
\bottomrule
\end{longtable}

\subsection{Key Directory Contents}
\paragraph{Project\_Task}
\begin{itemize}
  \item \texttt{RES201 Project.pdf}: official project brief and weekly deliverables.
  \item \texttt{ALIGNN\_Tutorial.pdf}: usage-oriented tutorial/reference for ALIGNN.
  \item \texttt{Project\_Timeline\_So\_Far.pdf}: human-written status summary from before this audit.
\end{itemize}

\paragraph{configs}
\begin{itemize}
  \item \texttt{oxide\_week1\_N50\_seed0.json}: earlier experimental config draft.
  \item \texttt{oxide\_week1\_N50\_seed0.fixed.json}: schema-compatible config produced during earlier debugging.
  \item \texttt{oxide\_week1\_N50\_seed0.run.json}: another historical run config.
  \item \texttt{oxide\_week1\_N50\_seed0.finetune\_last2.json}: final working oxide Week 1 config.
  \item \texttt{nitride\_week1\_N50\_seed0.finetune\_last2.json}: final working nitride Week 1 config.
\end{itemize}

\paragraph{docs}
\begin{itemize}
  \item \texttt{STAGE2\_PLAYBOOK.md}: Stage 2 instructions.
  \item \texttt{STAGE2\_FROZEN\_DATASET.md}: frozen dataset note file.
  \item \texttt{STAGE3\_PLAYBOOK.md}: Week 1 training output expectations.
\end{itemize}

\paragraph{env}
\begin{itemize}
  \item \texttt{bootstrap\_res201\_stage2\_data.sh}: create the Stage 2 data environment.
  \item \texttt{verify\_res201\_stage2\_data.py}: verify Stage 2 environment readiness.
  \item \texttt{bootstrap\_res201\_stage3\_train.sh}: create the Stage 3 training environment.
\end{itemize}

\paragraph{manifests}
\begin{itemize}
  \item \texttt{dft\_3d\_formation\_energy\_peratom\_splits.csv}: official split mapping used for this project.
  \item \texttt{dft\_3d\_formation\_energy\_peratom\_splits\_conflicts.csv}: ambiguous IDs excluded conservatively.
  \item \texttt{original\_split\_manifest\_template.csv}: template for user-supplied split manifests.
  \item \texttt{original\_split\_manifest\_template.json}: JSON split template.
\end{itemize}

\paragraph{scripts/dataset}
\begin{itemize}
  \item \texttt{build\_res201\_family\_datasets.py}: main Stage 2 dataset builder.
  \item \texttt{make\_split\_manifest\_from\_benchmark.py}: helper for constructing split manifests.
  \item \texttt{materialize\_alignn\_root.py}: materialize ALIGNN-ready directories.
  \item \texttt{res201\_stage2\_lib.py}: shared Stage 2 library logic.
  \item \texttt{validate\_res201\_stage2.py}: Stage 2 validator.
\end{itemize}

\paragraph{scripts/shared}
\begin{itemize}
  \item \texttt{zero\_shot\_alignn.py}: zero-shot evaluator.
  \item \texttt{prepare\_week1\_finetune\_dataset.py}: construct local $N=50$ experiment roots.
  \item \texttt{write\_week1\_alignn\_config.py}: write TrainingConfig-compatible JSONs.
  \item \texttt{finetune\_last2\_alignn.py}: repaired fine-tuning wrapper for this ALIGNN environment.
  \item \texttt{save\_pretrained\_state\_dict.py}: helper for extracting pretrained model weights.
  \item \texttt{check\_week1\_status.sh}: end-to-end Week 1 verification script.
  \item \texttt{generate\_workspace\_inventory.py}: generate the machine-readable full workspace inventory.
\end{itemize}

\paragraph{data\_shared}
\begin{itemize}
  \item \texttt{diagnostics/}: global split manifest, schema report, and catalog diagnostics.
  \item \texttt{oxide/manifests/}: \texttt{all.csv}, \texttt{train.csv}, \texttt{val.csv}, \texttt{test.csv}, \texttt{pool.csv}.
  \item \texttt{oxide/summaries/summary.json}: saved oxide Stage 2 summary.
  \item \texttt{oxide/structures/}: one POSCAR file per oxide-family record.
  \item \texttt{oxide/alignn\_ready/}: pool and test exports for ALIGNN consumption.
  \item \texttt{nitride/manifests/}: \texttt{all.csv}, \texttt{train.csv}, \texttt{val.csv}, \texttt{test.csv}, \texttt{pool.csv}.
  \item \texttt{nitride/summaries/summary.json}: saved nitride Stage 2 summary.
  \item \texttt{nitride/structures/}: one POSCAR file per nitride-family record.
  \item \texttt{nitride/alignn\_ready/}: pool and test exports for ALIGNN consumption.
\end{itemize}

\paragraph{results}
\begin{itemize}
  \item \texttt{Results_Before_Correction/oxide/zero\_shot/predictions.csv}: oxide zero-shot predictions.
  \item \texttt{Results_Before_Correction/nitride/zero\_shot/predictions.csv}: nitride zero-shot predictions.
  \item \texttt{reports/zero\_shot/zero\_shot\_summary.csv}: single canonical zero-shot MAE summary table.
  \item \texttt{Results_Before_Correction/oxide/N50\_seed0/dataset\_root/}: local oxide Week 1 POSCAR subset, \texttt{id\_prop.csv}, and split manifest.
  \item \texttt{Results_Before_Correction/nitride/N50\_seed0/dataset\_root/}: local nitride Week 1 POSCAR subset, \texttt{id\_prop.csv}, and split manifest.
  \item \texttt{Results_Before_Correction/oxide/N50\_seed0/finetune\_last2/}: completed oxide Week 1 fine-tune outputs.
  \item \texttt{Results_Before_Correction/nitride/N50\_seed0/finetune\_last2/}: completed nitride Week 1 fine-tune outputs.
  \item \texttt{Results_Before_Correction/oxide/N50\_seed0/train\_alignn\_restart/}: older partial attempt from the original debugging path.
\end{itemize}

\subsection{Full Inventory File}
Because the workspace contains thousands of generated POSCAR files and many derived artifacts, a full machine-generated inventory should accompany this report rather than being expanded inline.
The generated inventory artifacts for this audit are:
\begin{itemize}
  \item \texttt{reports/workspace\_inventory.csv}
  \item \texttt{reports/workspace\_inventory\_summary.json}
\end{itemize}
They record every non-\texttt{.git} path in the workspace.
At the time of generation, the inventory contained:
\begin{itemize}
  \item \texttt{19,298} total paths;
  \item \texttt{19,243} files; and
  \item \texttt{55} directories.
\end{itemize}
The CSV provides one row per path with a path string, item kind, and short description so that the full workspace can be audited without manually expanding thousands of structure files in the PDF itself.

\section{Git and GitHub Recommendation}
The repository remote is:
\begin{itemize}
  \item \url{https://github.com/TheArchitect999/res201-alignn-domain-shift}
\end{itemize}

Yes, the repository should be updated, but \textbf{not by pushing everything blindly}.
The correct policy is:
\begin{itemize}
  \item \textbf{Push:} scripts, configs, manifests, docs, summaries, lightweight prediction CSVs, report sources, and verification scripts.
  \item \textbf{Do not push:} heavy structure dumps, LMDB caches, model checkpoints, cached downloads, zipped local payloads, and large generated POSCAR collections inside experiment roots.
\end{itemize}

To make this safer, the local \texttt{.gitignore} was tightened during this audit to exclude:
\begin{itemize}
  \item root-level LMDB caches;
  \item pretrained checkpoint binary;
  \item heavy result checkpoints and LMDB caches;
  \item local Stage 2 zip archive;
  \item copied POSCARs inside \texttt{Results_Before_Correction/*/dataset\_root}.
\end{itemize}

Recommended push workflow:
\begin{lstlisting}[language=bash]
git status
git add .gitignore
git add scripts/shared/Check_Week1_Baseline_Status.sh
git add scripts/shared/Fine_Tune_Last_Two_ALIGNN_Layers.py
git add scripts/shared/Write_Week1_ALIGNN_Config.py
git add configs/oxide_week1_N50_seed0.finetune_last2.json
git add configs/nitride_week1_N50_seed0.finetune_last2.json
git add reports/zero_shot/zero_shot_summary.csv
git add Results_Before_Correction/oxide/zero_shot/predictions.csv
git add Results_Before_Correction/nitride/zero_shot/predictions.csv
git add Results_Before_Correction/oxide/N50_seed0/finetune_last2/*.json
git add Results_Before_Correction/oxide/N50_seed0/finetune_last2/*.csv
git add Results_Before_Correction/nitride/N50_seed0/finetune_last2/*.json
git add Results_Before_Correction/nitride/N50_seed0/finetune_last2/*.csv
git add reports/week1_report.tex
git add reports/workspace_inventory_summary.json
git add reports/workspace_inventory.csv
git commit -m "Complete RES201 week 1 oxide and nitride experiments"
git push origin main
\end{lstlisting}

Before pushing, run \texttt{git status} and verify that no huge \texttt{.pt}, \texttt{.vasp}, LMDB cache, or archive files are staged.

\section{Conclusion}
The workspace now satisfies the Week 1 requirements for both the oxide and nitride arms.
The work completed in this repository is sufficient to claim Week 1 completion, provided the team reports the outcomes honestly:
\begin{itemize}
  \item Stage 2 family datasets are valid and frozen;
  \item zero-shot results are recorded for both families;
  \item one $N=50$ fine-tuning run has been completed for each family;
  \item all required evidence files exist and can be verified from the terminal.
\end{itemize}

The next scientific priority is not Week 1 completion, which is already done, but explaining why the current $N=50$ fine-tunes underperform zero-shot and extending the study across additional $N$ values and seeds.

\end{document}
